\chapter{Kesimpulan dan Saran}

%=========================================================%
%                  TULIS PENUTUP DI SINI                  %
% Jangan Lupa untuk Menghapus Contoh Tulisan di Bawah Ini %
% Jangan Lupa untuk Menghapus Contoh Tulisan di Bawah Ini %
%          (Termasuk Panduan di Setiap Subbabnya)         %
% Contoh Teks Dapat Di-comment dengan Tanda "%" di Depan  %
%   jika Ingin Mempertahankan Panduan di Dalam File Ini   %
%=========================================================%

%------------------------ IMBAUAN ------------------------%
% Dilarang keras memunculkan sitasi baru (seperti Menurut %
% Gustin (2026)...). Bab ini adalah paparan milik penulis %
% sehingga tidak boleh ada pemikiran orang lain yang      %
% diselipkan di sini.                                     %
%---------------------------------------------------------%

% --- Kesimpulan ---
\section{Kesimpulan}

Sajikan resume akhir yang padat, jelas, dan lugas sebagai jawaban langsung atas seluruh butir Perumusan Masalah yang telah Anda ajukan di Bab I. Sebaiknya dijelaskan dengan pernyataan kualitatif mengenai makna dari hasil pengujian sehingga tidak berkesan menyalin ulang angka-angka statistik atau data mentah dari Bab IV. Selain itu, tegaskan kembali secara singkat kebaruan (\textit{novelty}) atau kontribusi utama yang berhasil ditemukan melalui penelitian ini. Tuliskan kesimpulan secara naratif berparagraf atau menggunakan penomoran (\textit{enumerate}) yang berkorespondensi satu-satu dengan jumlah rumusan masalah.



% --- Saran ---
\section{Saran}

Uraikan rekomendasi atau tindakan nyata yang dirumuskan secara logis berdasarkan poin-poin Kesimpulan di atas. Bagian ini tidak boleh berisi saran normatif yang bersifat umum (seperti ``semoga bermanfaat'' atau ``perlu ditingkatkan lagi''), melainkan harus spesifik dan operasional:

\begin{enumerate}
	\item \textbf{Saran Praktis/Aplikatif:} Ditujukan langsung kepada institusi, perusahaan, pelaku industri, atau masyarakat yang diteliti mengenai langkah konkret apa yang perlu mereka lakukan berdasarkan hasil temuan Anda.
	
	\item \textbf{Saran Akademis/Metodologis:} Ditujukan kepada peneliti selanjutnya yang tertarik untuk mengembangkan topik sejenis (misalnya keterbatasan instrumen yang perlu diperbaiki, variabel baru yang penting untuk ditambahkan, atau penggunaan metode lain untuk memperluas cakupan riset).
\end{enumerate}